\section{Catégories test et arche De Morgan}\label{sec:test}

Cette section reprend le ticket~\#187 de l'édition: les questions que les dossiers 108, 118 et 134-2 laissent ouvertes sur les catégories test, lues comme des ponts. On fixe un localisateur fondamental $W$ sur la catégorie $\mathrm{Cat}$ des petites catégories; les équivalences faibles usuelles, celles dont le nerf est une équivalence faible d'ensembles simpliciaux, forment le plus petit. Pour une petite catégorie $A$, $i_A\colon\widehat A\to\mathrm{Cat}$ envoie un préfaisceau $X$ sur sa catégorie des éléments $A/X$, et $i_A^*(C)(a) = \Hom(A/a, C)$. On suit les définitions de \cite[\S4.1]{cisinski-308}, qui reprennent \emph{À la poursuite des champs}: $A$ est une \emph{catégorie test faible} si $i_A^*(C)$ est $W$-asphérique pour toute petite catégorie $C$ à objet final (déf.~4.1.8); une \emph{catégorie test locale} si chaque $A/a$ est une catégorie test faible (déf.~4.1.12); une \emph{catégorie test} si elle est les deux. L'objet de Lawvere $L_A = i_A^*(\Delta_1)$ est le préfaisceau des cribles, et le critère de Grothendieck s'énonce: $A$ est une catégorie test locale si et seulement si $L_A$ est localement asphérique, c'est-à-dire si $A/(a\times L_A)$ est asphérique pour tout objet $a$ (\cite[th.~4.1.19]{cisinski-308}, \cite[th.~6.6]{maltsiniotis-2001}). C'est un pont au sens du \S\ref{sec:recall}, avec une arche homotopique: une propriété du topos $\widehat A$, la contractibilité de son objet des valeurs de vérité, lue sur la catégorie $A$.

\subsection{Les questions et ce qu'on en sait}

Le tableau donne les questions sur les catégories test, test faibles, locales et strictes, l'asphéricité et les intervalles séparants que les trois dossiers laissent ouvertes, par dossier et page de la transcription. Le dossier 134-2 ne contient de \emph{Pursuing Stacks} que les plans et tables des matières (le texte est imprimé); on y donne la section qu'annonce la table, lue dans l'édition arXiv \cite{ps-arxiv}. \q{Nôtre} marque une étape qui est de l'édition; \q{ouvert} veut dire: non trouvé dans les sources lues, qui sont \cite{cisinski-308,maltsiniotis-2001,ps-arxiv,hubert-thesis,hubert-dk}. Nous n'avons consulté ni les articles de Jardine, ni \emph{Higher categories and homotopical algebra} de Cisinski, ni l'Astérisque~301 imprimé.

\begin{center}\scriptsize\emergencystretch=3em
\begin{tabularx}{\linewidth}{>{\raggedright\arraybackslash}p{1.25cm}>{\raggedright\arraybackslash}X>{\raggedright\arraybackslash}p{5.2cm}}
\toprule
dossier, page & question & statut\\
\midrule
108, p.~2, q.~2 & $W$ et $W_A$ sont-ils saturés? & réglé: tout localisateur fondamental est fortement saturé \cite[prop.~4.2.4]{cisinski-308}; $W_A$ aussi (118, pp.~39--40)\\
108, p.~2, q.~3 & $(\mathrm{Cat})$ est-elle une catégorie de modèles? & réglé: structure de Thomason \cite[th.~5.2.12, 5.2.15]{cisinski-308}, pour tout $W$ accessible\\
108, p.~2, q.~4 & exactitude de $\mathrm{Hot}_W$ et de $\mathrm{Cat}\to\mathrm{Hot}_W$, $\widehat A\to\mathrm{Hot}_W$ & en partie: sommes et produits finis \cite[\S9]{maltsiniotis-2001}, limites inductives filtrantes \cite[cor.~4.2.22]{cisinski-308}, produits et $\widehat A\to\mathrm{Hot}$ ssi $A$ totalement asphérique \cite[prop.~4.3.2]{cisinski-308} (avec la correction de la ligne \nolinkurl{118-groupoid-products-not-totally-aspherical})\\
108, p.~2, q.~5 & quand $\widehat A$ est-elle une catégorie de modèles fermée? & réglé: une structure à cofibrations les monomorphismes et équivalences faibles $W_A$ existe ssi $A$ est test locale \cite[introduction, th.~1.4.3, cor.~4.2.18]{cisinski-308}\\
108, p.~2, q.~6 & $\Delta^{\circ}$ est-elle une catégorie test? (\q{non}) & réglé, nôtre: $\Delta^\circ$ a un objet initial, donc n'est ni test faible ni test locale (théorème~\ref{th:ore}); le cas test est énoncé sans argument dans \cite[\S88]{ps-arxiv}\\
108, p.~2, q.~6a & les variantes $\Delta_f$, cubes & $\Delta_f$ est test faible, non test \cite[\S8]{maltsiniotis-2001}; la catégorie des cubes est test \cite[cor.~8.4.13]{cisinski-308}; les autres signes cubiques ne sont pas identifiés sur la page\\
108, p.~2, q.~7, 17 & invariants de $\widehat A$, Dold--Puppe & en partie: correspondances de Dold--Kan homotopiques pour les catégories test vérifiant une condition de Whitehead, $\Theta$ comprise \cite{hubert-thesis,hubert-dk}\\
108, p.~3, q.~13 & asphérique à topos 0-connexe $\Rightarrow$ totalement asphérique? & en partie: l'asphéricité totale est celle des $a\times b$ \cite[4.3.1--4.3.2]{cisinski-308}, \cite[prop.~7.1]{maltsiniotis-2001}; une catégorie test faible totalement asphérique est test stricte \cite[prop.~1.17]{hubert-dk}; la forme 0-connexe de la page: ouvert\\
108, p.~4, q.~16 & pour un contracteur $A$, $i_A$, $j_A$ inverses à homotopismes près; contractibilité par $A/X$; $i_A(X\times Y)\to i_AX\times i_AY$? & à équivalences faibles près: oui, un contracteur non trivial étant test strict (\cite[\S88]{ps-arxiv}, \cite[\S7]{maltsiniotis-2001}) et \cite[th.~4.1.26, prop.~4.3.2]{cisinski-308}; à homotopismes près: ouvert\\
118, p.~14, pp.~18--22 & sommes amalgamées le long d'immersions dans $(\mathrm{Cat})$ & la page répond \q{non} p.~22; critères de Cisinski \cite[\S5.1]{cisinski-308} (ligne \nolinkurl{118-pushout-along-cofibrations-in-cat})\\
118, pp.~23--25 & calcul des fractions pour $\mathrm{Hot}_A$ avec un intervalle séparant; quelle classe $\mathfrak M$ dans $(\mathrm{Cat})$? & dans $\widehat A$ la page demande $A$ test locale, comme \cite[th.~4.1.19]{cisinski-308}; la version dans $(\mathrm{Cat})$: ouvert\\
118, p.~14 (1) & $W = W^{\circ}$? & réglé \cite[sorite~3.3.8 (7)]{cisinski-308}\\
118, pp.~26--30 & sommes quelconques dans $\mathrm{Hot}(W)$ & réglé pour $W$ accessible, nôtre à partir de \cite[3.3.8 (1), th.~5.2.15]{cisinski-308}\\
118, p.~36 & un produit infini d'objets asphériques est-il asphérique? & non, nôtre (proposition~\ref{prop:products})\\
134-2, p.~5, \S\S31--33 & un critère pour les catégories test, nécessité de T1--T3 & réglé \cite[th.~4.1.19]{cisinski-308}, \cite[th.~6.6]{maltsiniotis-2001}\\
134-2, p.~6 \S43, p.~9 \S98 & $\Delta^f$ test faible; \q{perplexities about its being aspheric} & réglé \cite[\S8]{maltsiniotis-2001}\\
134-2, p.~9, \S84 & l'inclusion $f_!(M_{\mathrm{as}})\subset M'_{\mathrm{as}}$ & laissée sans réponse dans \cite[\S88, 1)]{ps-arxiv}; ouvert\\
134-2, p.~9, \S86 & existe-t-il une notion de morphisme de structures d'asphéricité? & ouvert\\
134-2, p.~9, \S90 & existence de foncteurs test (parfaits) vers un modélisateur & ouvert en général; pour $\Delta\to\mathrm{Cat}$, le $c\,\mathrm{Sd}^2$ de Thomason donne une équivalence de Quillen \cite[th.~5.2.12]{cisinski-308}\\
134-2, p.~9, \S91 & conditions sur $A$ pour les types d'homotopie de type fini; catégories test à opérations de bord & ouvert; pour l'homologie, \cite{hubert-dk}\\
\bottomrule
\end{tabularx}
\end{center}

\subsection{L'obstruction De Morgan}

Par l'arche De Morgan (arche~B de \cite{bridges-generator}), $\widehat A$ est De Morgan si et seulement si $A$ vérifie la condition d'Ore. L'objet de Lawvere voit directement la condition d'Ore.

\begin{lemma}\label{lem:ore-sieves}
$A$ vérifie la condition d'Ore si et seulement si, pour toute flèche $f\colon b\to a$ et tout crible non vide $S$ sur $a$, le crible $f^*S$ est non vide. Autrement dit, les cribles vides forment un facteur direct de $L_A$: $L_A = L_A^{\emptyset}\sqcup L_A^{+}$, où $L_A^{+}(a)$ est l'ensemble des cribles non vides sur $a$.
\end{lemma}

\begin{proof}
Si $S$ contient $g\colon c\to a$ et si le cospan $(f,g)$ se complète en $fh = gk$, alors $h\in f^*S$. Inversement, étant donné un cospan $(f,g)$, on applique la condition au crible engendré par $g$: une flèche $h$ de $f^*S$ est un $h$ tel que $fh = gk$ pour un $k$. Les cribles vides forment toujours un sous-préfaisceau; leur complémentaire en est un si et seulement si aucun crible non vide ne se tire en le crible vide.
\end{proof}

\begin{theorem}\label{th:ore}
Soient $W$ un localisateur fondamental autre que les deux localisateurs grossiers, et $A$ une petite catégorie non vide vérifiant la condition d'Ore, c'est-à-dire telle que $\widehat A$ soit De Morgan. Alors $A$ n'est ni une $W$-catégorie test faible ni une $W$-catégorie test locale. Plus généralement, s'il existe un préfaisceau $X$ tel que $A/X$ vérifie la condition d'Ore, $A$ n'est pas une $W$-catégorie test locale.
\end{theorem}

\begin{proof}
Tout localisateur fondamental non grossier est contenu dans le localisateur $W_0$ des foncteurs qui induisent une bijection sur les $\pi_0$ \cite[prop.~9.3.2]{cisinski-308}; une catégorie $W$-asphérique est donc connexe. Par le lemme~\ref{lem:ore-sieves}, $A/L_A$ est la réunion disjointe de $A/L_A^{\emptyset}\simeq A$ et de $A/L_A^{+}$, qui contient les cribles maximaux; toutes deux sont non vides, donc $L_A$ n'est pas asphérique, alors qu'une catégorie test faible l'exige (déf.~4.1.8 avec $C=\Delta_1$). De même $a\times L_A = (a\times L_A^{\emptyset})\sqcup(a\times L_A^{+})$, chaque morceau ayant une section au-dessus de $a$, donc $L_A$ n'est pas localement asphérique et $A$ n'est pas test locale \cite[th.~4.1.19]{cisinski-308}. La dernière assertion en résulte, puisque $A/X$ est test locale dès que $A$ l'est \cite[prop.~4.1.14 (v)]{cisinski-308}, et qu'une tranche d'une catégorie non vide est non vide.
\end{proof}

\begin{corollary}\label{cor:ore}
Pour les équivalences faibles usuelles, aucune des catégories suivantes n'est test faible ni test locale: un groupoïde non vide; un monoïde commutatif, un groupe, un monoïde à zéro; un monoïde fini à un seul idéal à droite minimal; une catégorie à objet initial, comme $\Delta^{\circ}$, la catégorie des simplexes augmentée, ou les ensembles finis avec l'ensemble vide; la catégorie libre sur un carquois dont chaque sommet reçoit au plus une flèche (\lean{rightOreCondition_paths_iff}). Aucune des suivantes n'est test locale: un ensemble ordonné à élément minimal, la tranche au-dessus de cet élément étant le point; une catégorie libre finie sur un carquois sans cycle, pour la même raison en une source.
\end{corollary}

Pour un monoïde fini, la condition d'Ore est l'existence d'un seul idéal à droite minimal: deux idéaux à droite minimaux sont disjoints, et tout idéal à droite $mM$ en contient un. Pour la seconde liste on applique la dernière assertion du théorème~\ref{th:ore} avec $X$ représentable.

\status{} Grothendieck énonce le cas test pour un objet initial dans \emph{Pursuing Stacks}: \q{it is easy to see that a test category cannot possibly have an initial object} \cite[\S88, p.~294 du tapuscrit]{ps-arxiv}, sans argument; Maltsiniotis remarque qu'un segment séparant représentable exclut un objet initial \cite[rem.~6.8]{maltsiniotis-2001}, et les familles de catégories test de \cite[ex.~4.1.20]{cisinski-308} et \cite[ex.~6.9]{maltsiniotis-2001} excluent la catégorie vide. La réponse \q{non} du dossier 108 à la question~6 s'accorde donc avec une remarque ultérieure de Grothendieck lui-même. Ce que nous n'avons trouvé ni dans ces sources, ni par TheoremSearch et alphaXiv, c'est la forme générale: la condition d'Ore, c'est-à-dire la propriété De Morgan du topos, exclut les catégories test faibles et test locales. Le pont la lit en une ligne: une propriété logique de $\widehat A$ (la loi de De Morgan) interdit une propriété homotopique (un objet des valeurs de vérité contractile). La réciproque est fausse: $\Delta$ ne vérifie pas la condition d'Ore et est une catégorie test, tandis que les monoïdes finis non Ore ci-dessous n'en sont pas.

\subsection{Ce que les arches engendrent}

Composé avec les arches déjà démontrées, le théorème~\ref{th:ore} donne trois familles d'énoncés.

\emph{Tabliers d'Ore sans tranche ponctuelle.} L'argument élémentaire du corollaire~\ref{cor:ore} demande un objet qui ne reçoit que son identité; la condition d'Ore n'en demande pas. Par la forme préordonnée de l'arche De Morgan (\lean{rightOreCondition_preorder_iff}), un préordre non vide dans lequel chaque $\downarrow x$ est filtrant vers le bas n'est ni test faible ni test local: c'est le cas de $\mathbb Z$, $\mathbb Q$, $\mathbb R$ et de tout ordre total, bien qu'aucune de leurs tranches ne soit le point. Il en va de même du monoïde $\mathbb N$, dont les préfaisceaux sont les ensembles munis d'un endomorphisme (il n'est pas test local, et $B\mathbb N\simeq S^1$ exclut déjà qu'il soit test), et de la catégorie libre sur un carquois dont chaque sommet reçoit une flèche et aucun n'est une source, comme $\cdots\to-1\to0\to1\to\cdots$. Parmi les 2510 monoïdes d'ordre au plus~6 de l'atlas, 2015 vérifient la condition d'Ore et relèvent du théorème sans calcul d'homologie.

\emph{Les catégories test ne sont pas De Morgan.} Lu à l'envers, le théorème~\ref{th:ore} dit que le topos de préfaisceaux d'une catégorie test faible ou test locale n'est jamais De Morgan. Pour les ensembles simpliciaux, c'était démontré directement (\lean{not_isDeMorganTopos_simplicialSets}); cela s'ensuit désormais pour les ensembles cubiques \cite[cor.~8.4.13]{cisinski-308}, pour les $\Theta_n$-ensembles et les $\Theta$-ensembles \cite[th.~6.5]{hubert-dk}, et pour toutes les catégories test rappelées dans \cite[introduction]{hubert-dk}, la catégorie dendroïdale comprise.

\begin{proposition}\label{prop:slices}
Soient $W$ un localisateur fondamental non grossier et $A$ une $W$-catégorie test locale. Pour tout préfaisceau non vide $X$ sur $A$, le topos tranche $\widehat A/X\simeq\widehat{A/X}$ n'est pas De Morgan. En particulier, pour tout objet $a$, un objet muni d'une flèche vers $a$ est le but d'un cospan qui ne se complète pas en carré commutatif.
\end{proposition}

\begin{proof}
$A/X$ est une catégorie test locale \cite[prop.~4.1.14 (v)]{cisinski-308}, non vide puisque $X$ l'est; par le théorème~\ref{th:ore}, elle ne vérifie pas la condition d'Ore, et par l'arche De Morgan $\widehat{A/X}$ n'est pas De Morgan. Pour $X = a$, un cospan de $A/a$ est un cospan de $A$ dont le sommet s'envoie sur $a$, et un carré commutatif dans $A/a$ en est un dans $A$.
\end{proof}

Un topos test n'est donc De Morgan nulle part, et pas seulement globalement. Nous n'avons pas cherché la proposition~\ref{prop:slices} dans la littérature; elle découle en deux lignes du théorème et de la proposition~4.1.14 de Cisinski.

\subsection{Tabliers finis: calculs et conjecture}

Sur les tabliers qu'énumère le générateur, le script \nolinkurl{computations/187-test-categories/finite_test.py} calcule l'homologie en degrés $0$ à $2$, sur $\mathbb Q$ et $\mathbb F_2$, des nerfs de $A$, de $A/L_A$ et de chaque $A/(a\times L_A)$. Pour tout monoïde d'ordre au plus~6 (2510 à isomorphisme près, d'après l'atlas de l'article associé), l'homologie de $A/L_A$ en degrés au plus~2 diffère de celle d'un point, et celle de $A/(a\times L_A)$ aussi: aucun monoïde d'ordre $\le 6$ n'est test faible ni test local, pour les équivalences faibles usuelles. Les monoïdes non Ore, au nombre de 495, échappent au théorème~\ref{th:ore}; pour eux $H_1(A/(a\times L_A);\mathbb Q)$ est de rang compris entre $4$ et $743$. Pour les ensembles ordonnés, \nolinkurl{posets_euler.py} calcule la caractéristique d'Euler de l'ensemble ordonné fini $A/L_A$ par la formule $\chi(A/L_A) = \sum_x |D(\downarrow x)|\,c(x)$, où $D(\downarrow x)$ est l'ensemble des parties commençantes de $\downarrow x$ et $c(x) = 1-\chi(A_{<x})$: pour tout ensemble ordonné à au plus six éléments et de caractéristique $\chi(A)=1$, $\chi(A/L_A)\ne1$, donc aucun ensemble ordonné à au plus six éléments n'est test faible (un ensemble ordonné de caractéristique $\chi(A)\ne1$ n'est même pas asphérique).

\begin{remark}[Conjecture]
Aucune catégorie finie non vide n'est test faible, ni test locale. Le cas local est établi pour les ensembles ordonnés finis et les catégories libres finies sans cycle (corollaire~\ref{cor:ore}); pour les monoïdes nous n'avons que les calculs ci-dessus.
\end{remark}

\subsection{Produits infinis}

\begin{proposition}\label{prop:products}
Pour tout localisateur fondamental non grossier $W$, un produit infini de catégories $W$-asphériques n'est pas nécessairement $W$-asphérique, et la classe $W$ n'est pas stable par produits infinis.
\end{proposition}

\begin{proof}
Soit $Z_n$ le zigzag $0\to1\leftarrow2\to\cdots$ à $n+1$ objets, un ensemble ordonné contractile. Un zigzag de longueur $k$ dans $\prod_n Z_n$ se projette sur des zigzags de longueur $k$ dans chaque facteur, et les extrémités de $Z_n$ sont à distance $n$; donc $(0)_n$ et $(n)_n$ sont dans deux composantes de $\prod_n Z_n$, qui n'est pas $W$-asphérique (démonstration du théorème~\ref{th:ore}). Ainsi $\prod_n(Z_n\to e)$ n'est pas dans $W$.
\end{proof}

C'est la réponse à la question du dossier 118, p.~36, que Grothendieck y juge moins importante que celle des sommes infinies; l'argument est celui de la p.~28, où les mêmes zigzags montrent que le passage $(\mathrm{Cat})\to\overline{\mathrm{Cat}}$ aux classes d'homotopie de foncteurs ne commute pas aux sommes infinies. Comme le nerf commute aux limites projectives, on obtient aussi un produit infini d'ensembles simpliciaux asphériques qui n'est pas asphérique.

\leanentry{} Six fichiers de \texttt{lean/bridges}. \lean{TestCategories.lean} démontre le noyau de connexité avec l'objet de Lawvere de Mathlib, \lean{Functor.sieves}: le lemme~\ref{lem:ore-sieves} (\lean{rightOreCondition_iff_pullback_ne_bot}), la décomposition de toute catégorie des éléments au-dessus de $L_A$ (\lean{not_isPreconnected_elements}) et ses conséquences (\lean{not_isConnected_elements_sieves}, \lean{not_isConnected_elements_yonedaProdSieves}, \lean{not_isConnected_elements_sieves_of_isDeMorganTopos}, \lean{not_isDeMorganTopos_of_isConnected_elements_sieves}). Mathlib n'a ni les équivalences faibles de $\mathrm{Cat}$ ni les notions test de Grothendieck, et la bibliothèque Stacks d'AxiomMath, que nous avons examinée, non plus; \lean{FundamentalLocalizers.lean} les pose donc axiomatiquement, comme Grothendieck: un localisateur fondamental (\lean{FundamentalLocalizer}: saturation faible, asphéricité des catégories à objet final, théorème~A relatif), $i_A^*$ (\lean{iStar}), les catégories test faibles et locales (\lean{IsWeakTest}, \lean{IsLocalTest}, d'après les déf.~4.1.8 et~4.1.12 de \cite{cisinski-308}). Sous l'hypothèse que les catégories asphériques sont connexes (\lean{AsphericalConnected}, ce que donne $W\subset W_0$), il démontre le théorème~\ref{th:ore} en entier (\lean{not_isWeakTest}, \lean{not_isLocalTest}, \lean{not_isLocalTest_of_over}), sa lecture par l'arche De Morgan et à l'envers (\lean{not_isWeakTest_of_isDeMorganTopos}, \lean{not_isDeMorganTopos_of_isWeakTest}), la proposition~\ref{prop:slices} pour tout préfaisceau $X$ qui a une section (\lean{not_isDeMorganTopos_slice_of_isLocalTest}, d'abord pour $X$ représentable, \lean{not_isDeMorganTopos_over_of_isLocalTest}), sans la proposition~4.1.14 de Cisinski: si $A/X$ vérifiait la condition d'Ore, sa tranche au-dessus d'un objet $p=(a,x)$ la vérifierait aussi, puis, le long d'une section à isomorphisme près de la projection $(A/X)/p\to A/a$ (\lean{overCostructuredArrowSection}, \lean{rightOreCondition_of_iso}), la catégorie test faible $A/a$, et pour le topos tranche $\widehat A/X$ lui-même par l'équivalence $\widehat A/X\simeq\widehat{A/X}$ de Mathlib (\lean{overEquivPresheafCostructuredArrow}; \lean{isDeMorganTopos_over_iff}, \lean{not_isDeMorganTopos_overPresheaf_of_isLocalTest}), les corollaires pour $\Delta^{\circ}$, les groupoïdes, les préordres filtrants et $\mathbb Z$, le monoïde $\mathbb N$, les monoïdes finis à un seul idéal à droite minimal, les catégories libres à au plus une flèche entrante par sommet, les préordres à élément minimal, et les catégories libres finies sans cycle, par l'existence d'une source (\lean{not_isTest_simplexCategory_op}, \lean{not_isTest_groupoid}, \lean{not_isTest_preorder}, \lean{not_isTest_int}, \lean{not_isTest_multiplicative_nat}, \lean{not_isTest_finite_monoid}, \lean{not_isTest_paths}, \lean{not_isLocalTest_of_minimal}, \lean{exists_source_of_acyclic}, \lean{not_isLocalTest_paths_of_acyclic}), et la proposition~\ref{prop:products} en entier: chaque zigzag $Z_n$ est asphérique pour tout localisateur fondamental, à partir des seuls axiomes (\lean{isAspherical_zigZag}: récurrence sur $n$, par le théorème~A relatif appliqué à l'inclusion $Z_n\to Z_{n+1}$ ou à la rétraction $Z_{n+1}\to Z_n$, dont les tranches ont un objet final), tandis que leur produit n'est pas connexe, donc pas asphérique (\lean{not_isPreconnected_zigZagProd}, \lean{isAspherical_zigZag_and_not_isAspherical_zigZagProd}). Non formalisés: que les équivalences faibles usuelles, définies par le nerf, forment un localisateur fondamental contenu dans $W_0$, et la proposition~9.3.2 de Cisinski (il faudrait les équivalences faibles d'ensembles simpliciaux); et le statut de catégorie test de $\Theta$ et de $\Omega$. \lean{Theta2.lean} définit la catégorie $\Theta_2=\Delta\wr\Delta$ de Joyal, qu'aucune bibliothèque n'a (\lean{Theta2}: un morphisme est une application de $\Delta$ sur les sommets avec, pour chaque cellule du but au-dessus d'une cellule de la source, une application de $\Delta$ entre leurs $2$-cellules), et démontre qu'elle est asphérique et ne vérifie pas la condition d'Ore (\lean{Theta2.isAspherical}, \lean{Theta2.not_rightOreCondition}): l'obstruction du théorème~\ref{th:ore} ne s'y applique donc pas. Son statut de catégorie test n'est pas formalisé. La méthode du cône échoue: une $1$-cellule finale demanderait un sommet dans chaque $[b_j]$ au-delà de $\varphi(m)$, et aucun choix n'est naturel. La méthode du cylindre échoue aussi: épaissir les $2$-cellules donne un foncteur à deux bouts naturels, mais une flèche qui écrase une cellule oublie les composantes au-dessus d'elle, et la contraction n'est pas naturelle. Les preuves connues passent par le critère de Grothendieck \cite[th.~4.1.19]{cisinski-308} ou par la stabilité des catégories test par $\Delta\wr-$. \lean{GrothendieckCriterion.lean} démontre ce critère à partir des seuls axiomes: $A$ est une catégorie test locale si et seulement si $L_A=i_A^*(\Delta_1)$ est localement asphérique (\lean{isLocalTest_iff_isLocallyAspherical}). Un sens se réduit au cas $C=\Delta_1$, en comparant $(A/a)/i_{A/a}^*(C)$ et $A/(a\times i_A^*C)$ (\lean{isAspherical_slice_iff}). L'autre est l'argument d'homotopie de Grothendieck. Quand $L_A$ est localement asphérique, la projection $A/(a\times i_A^*(\Delta_1\times C))\to A/(a\times i_A^*C)$ est une $W$-équivalence: ses tranches se rétractent sur $A/(b\times L_A)$, par le théorème~A (\lean{pmap_mem}). La contraction $\Delta_1\times C\to C$ sur l'objet final rétracte alors $A/(a\times i_A^*C)$ sur $A/a$ (\lean{isAspherical_prodY_of_pmap}). Pour $\Theta_2$, il reste à montrer que $L_{\Theta_2}$ est localement asphérique. \lean{SimplexTest.lean} démontre, à partir des seuls axiomes et sans hypothèse de connexité, que $\Delta$ est une catégorie test faible et locale pour tout localisateur fondamental (\lean{isWeakTest_simplexCategory}, \lean{isLocalTest_simplexCategory}, \lean{isTest_simplexCategory}). Deux lemmes portent l'homotopie: une transformation naturelle entre deux foncteurs fait de l'un une $W$-équivalence si et seulement si l'autre en est une (\lean{mem_iff_of_natTrans}, par la projection $\Delta_1\times D\to D$ et le théorème~A relatif), et une catégorie munie d'un endofoncteur $t$ et de transformations naturelles $1\to t\leftarrow\mathrm{const}_d$ est asphérique (\lean{isAspherical_of_contraction}). Pour $C$ à objet final $\star$, $A/i_A^*(C)$ se contracte par $(a,F)\mapsto(Ta,F^+)$, où $Ta$ ajoute à $a$ un sommet final et où $F^+$ vaut $F$ sur les simplexes de $Ta$ qui évitent le nouveau sommet et $\star$ sur les autres (\lean{ConeData}, \lean{ConeData.isWeakTest}); $\Delta$ et chaque tranche $\Delta/b$ portent cette structure de cône (\lean{coneData}, \lean{coneDataOver}, le nouveau sommet étant au-dessus du sommet final de $b$). \lean{CubeTest.lean} fait de même pour la catégorie des cubes $\square$ (faces et dégénérescences, \lean{Cube}; \lean{isTest_cube}, \cite[cor.~8.4.13]{cisinski-308} pour les équivalences faibles usuelles). $\square$ n'a pas de sommet de cône naturel, mais il a un cylindre $[1]^n\mapsto[1]^{n+1}$ à deux bouts; le bout $1$ donne une transformation naturelle de $(a,F)\mapsto(a,\star)$ vers la contraction, si bien que la projection $A/i_A^*(C)\to A$, rétraction de $a\mapsto(a,\star)$, est une $W$-équivalence, et $A/i_A^*(C)$ est asphérique puisque $A$ a un objet final (\lean{CylData}, \lean{CylData.isWeakTest}); chaque tranche $\square/b$ porte le cylindre au-dessus de $b$ par la projection (\lean{cylDataOver}).
